\section{Discussion of the Extensions}\label{sec:extension_limits}
Recovery tests show improved functionality on the selected difficult
inputs, with a larger gain on x64. The full-corpus wrapping and LLM
campaigns complement this comparison by assessing the added modes on all
102 inputs.

The controlled suites help explain what each stage contributes. All six
annotated programs are classified exactly, while the 13 functional
programs check the output of individual programming patterns. Their small
size and simple layouts make them useful diagnostic tests. Their high
success rates are not a consequence of being benign: the original corpus
contains larger and more varied instruction/data layouts, which exercise
additional cases. The recorded native x86 IBP result remains associated with its tested
version; a future corrected implementation requires a separately identified
evaluation.

The two LLM roles remain distinct. Data identification feeds shared
handling, giving identical binaries in 16/19 original cases and matching
17/19 functional outcomes. Assembly transformation proposes a new listing;
zero-shot and two-shot have similar aggregate outcomes. Input size,
construction success, and execution outcome offer more informative
distinctions in this campaign.

The empirical scope can be stated once. Functional preservation is assessed
through the specified execution tests, rather than a formal proof for every
possible execution. Classification accuracy is measured on independently
annotated programs; the unannotated payload comparison measures agreement
and downstream functionality. Static detection rates describe the submitted
raw files and the reported engine responses, including the disclosed
unsupported and failed analyses. Repeated inputs, shared hashes, recorded
version differences, and retries are retained in the artifact so that these
descriptive results can be inspected and reproduced.
